% problem_id: S001-lemma-base-change-does-not-hold-pre
% source_id: S001 (Stacks Project)
% source_file: etale-cohomology.tex
% statement_locator: etale-cohomology.tex lines 17078--17094
% proof_locator: etale-cohomology.tex lines 17096--17293
% env: lemma
% id_note: label 在本来源内唯一
% pairing: tight (verified)
% status: complete (statement + author proof verbatim + reason + locators)
%
\begin{lemma}
\label{lemma-base-change-does-not-hold-pre}
With $f : X \to S$ and $n$ as in Remark \ref{remark-base-change-holds}
assume for some $q \geq 1$ we have that $BC(f, n, q - 1)$ is true, but
$BC(f, n, q)$ is not. Then there exist a commutative diagram
$$
\xymatrix{
X \ar[d]_f & X' \ar[d]_{f'} \ar[l] & Y \ar[l]^h \ar[d]^e \\
S & S' \ar[l] & \Spec(K) \ar[l]_g
}
$$
where $X' = X \times_S S'$, $Y = X' \times_{S'} \Spec(K)$,
$K$ is a field, and $\mathcal{F}$ is an abelian sheaf
on $\Spec(K)$ annihilated by $n$ such that
$(f')^{-1}R^qg_*\mathcal{F} \to R^qh_*e^{-1}\mathcal{F}$
is not an isomorphism.
\end{lemma}

\begin{proof}
Choose a commutative diagram
$$
\xymatrix{
X \ar[d]_f & X' \ar[l] \ar[d]_{f'} & Y \ar[l]^h \ar[d]^e \\
S & S' \ar[l] & T \ar[l]_g
}
$$
with $X' = X \times_S S'$ and $Y = X' \times_{S'} T$ and
$g$ quasi-compact and quasi-separated, and an abelian sheaf
$\mathcal{F}$ on $T_\etale$ annihilated by $n$ such that
the base change map
$(f')^{-1}R^qg_*\mathcal{F} \to R^qh_*e^{-1}\mathcal{F}$
is not an isomorphism. Of course we may and do replace $S'$
by an affine open of $S'$; this implies that $T$ is quasi-compact
and quasi-separated. By Lemma \ref{lemma-base-change-q-injective} we see
$(f')^{-1}R^qg_*\mathcal{F} \to R^qh_*e^{-1}\mathcal{F}$
is injective. Pick a geometric point $\overline{x}$ of $X'$
and an element $\xi$ of $(R^qh_*q^{-1}\mathcal{F})_{\overline{x}}$
which is not in the image of the map
$((f')^{-1}R^qg_*\mathcal{F})_{\overline{x}} \to
(R^qh_*e^{-1}\mathcal{F})_{\overline{x}}$.

\medskip\noindent
Consider a morphism $\pi : T' \to T$ with $T'$ quasi-compact
and quasi-separated and denote $\mathcal{F}' = \pi^{-1}\mathcal{F}$.
Denote $\pi' : Y' = Y \times_T T' \to Y$ the base change of $\pi$
and $e' : Y' \to T'$ the base change of $e$. Picture
$$
\vcenter{
\xymatrix{
X' \ar[d]_{f'} & Y \ar[l]^h \ar[d]^e & Y' \ar[l]^{\pi'} \ar[d]^{e'} \\
S' & T \ar[l]_g & T' \ar[l]_\pi
}
}
\quad\text{and}\quad
\vcenter{
\xymatrix{
X' \ar[d]_{f'} & & Y' \ar[ll]^{h' = h \circ \pi'} \ar[d]^{e'} \\
S' & & T' \ar[ll]_{g' = g \circ \pi}
}
}
$$
Using pullback maps we obtain a canonical commutative diagram
$$
\xymatrix{
(f')^{-1}R^qg_*\mathcal{F} \ar[r] \ar[d] &
(f')^{-1}R^qg'_*\mathcal{F}' \ar[d] \\
R^qh_*e^{-1}\mathcal{F} \ar[r] &
R^qh'_*(e')^{-1}\mathcal{F}'
}
$$
of abelian sheaves on $X'$. Let $P(T')$ be the property
\begin{itemize}
\item The image $\xi'$ of $\xi$ in
$(Rh'_*(e')^{-1}\mathcal{F}')_{\overline{x}}$ is
not in the image of the map
$(f^{-1}R^qg'_*\mathcal{F}')_{\overline{x}} \to
(R^qh'_*(e')^{-1}\mathcal{F}')_{\overline{x}}$.
\end{itemize}
We claim that hypotheses (1), (2), and (3) of
Lemma \ref{lemma-formal-argument} hold for $P$
which proves our lemma.

\medskip\noindent
Condition (1) of Lemma \ref{lemma-formal-argument}
holds for $P$ because the \'etale topology of a scheme
and a thickening of the scheme is the same. See
Proposition \ref{proposition-topological-invariance}.

\medskip\noindent
Suppose that $I$ is a directed set and that $T_i$
is an inverse system over $I$ of quasi-compact and quasi-separated
schemes over $T$ with affine transition morphisms.
Set $T' = \lim T_i$. Denote $\mathcal{F}'$ and $\mathcal{F}_i$
the pullback of $\mathcal{F}$ to $T'$, resp.\ $T_i$. Consider
the diagrams
$$
\vcenter{
\xymatrix{
X \ar[d]_{f'} & Y \ar[l]^h \ar[d]^e & Y_i \ar[l]^{\pi_i'} \ar[d]^{e_i} \\
S & T \ar[l]_g & T_i \ar[l]_{\pi_i}
}
}
\quad\text{and}\quad
\vcenter{
\xymatrix{
X \ar[d]_{f'} & & Y_i \ar[ll]^{h_i = h \circ \pi_i'} \ar[d]^{e_i} \\
S & & T_i \ar[ll]_{g_i = g \circ \pi_i}
}
}
$$
as in the previous paragraph. It is clear that $\mathcal{F}'$ on
$T'$ is the colimit of the pullbacks of $\mathcal{F}_i$ to $T'$
and that $(e')^{-1}\mathcal{F}'$ is the colimit of the pullbacks
of $e_i^{-1}\mathcal{F}_i$ to $Y'$.
By Lemma \ref{lemma-relative-colimit-general}
we have
$$
R^qh'_*(e')^{-1}\mathcal{F}' = \colim R^qh_{i, *}e_i^{-1}\mathcal{F}_i
\quad\text{and}\quad
(f')^{-1}R^qg'_*\mathcal{F}' = \colim (f')^{-1}R^qg_{i, *}\mathcal{F}_i
$$
It follows that if $P(T_i)$ is true for all $i$, then
$P(T')$ holds. Thus condition (2) of Lemma \ref{lemma-formal-argument}
holds for $P$.

\medskip\noindent
The most interesting is condition (3) of Lemma \ref{lemma-formal-argument}.
Assume $T'$ is a quasi-compact and quasi-separated scheme over $T$
such that $P(T')$ is true.
Let $Z \subset T'$ be a closed subscheme with complement $V \subset T'$
quasi-compact. Consider the diagram
$$
\xymatrix{
Y' \times_{T'} Z \ar[d]_{e_Z} \ar[r]_{i'} &
Y' \ar[d]_{e'} &
Y' \times_{T'} V \ar[l]^{j'} \ar[d]^{e_V} \\
Z \ar[r]^i &
T' &
V \ar[l]_j
}
$$
Choose an injective map $j^{-1}\mathcal{F}' \to \mathcal{J}$
where $\mathcal{J}$ is an injective sheaf of $\mathbf{Z}/n\mathbf{Z}$-modules
on $V$. Looking at stalks we see that the map
$$
\mathcal{F}' \to \mathcal{G} = j_*\mathcal{J} \oplus i_*i^{-1}\mathcal{F}'
$$
is injective. Thus $\xi'$ maps to a nonzero element of
\begin{align*}
& \Coker\left(
((f')^{-1}R^qg'_*\mathcal{G})_{\overline{x}}
\to
(R^qh'_*(e')^{-1}\mathcal{G})_{\overline{x}}
\right) = \\
&
\Coker\left(
((f')^{-1}R^qg'_*j_*\mathcal{J})_{\overline{x}}
\to
(R^qh'_*(e')^{-1}j_*\mathcal{J})_{\overline{x}}
\right) \oplus \\
& \Coker\left(
((f')^{-1}R^qg'_*i_*i^{-1}\mathcal{F}')_{\overline{x}}
\to
(R^qh'_*(e')^{-1}i_*i^{-1}\mathcal{F}')_{\overline{x}}
\right)
\end{align*}
by part (2) of Lemma \ref{lemma-base-change-q-injective}.
If $\xi'$ does not map to zero in the second summand, then
we use
$$
(f')^{-1}R^qg'_*i_*i^{-1}\mathcal{F}' =
(f')^{-1}R^q(g' \circ i)_*i^{-1}\mathcal{F}'
$$
(because $Ri_* = i_*$ by
Proposition \ref{proposition-finite-higher-direct-image-zero}) and
$$
R^qh'_*(e')^{-1}i_*i^{-1}\mathcal{F} =
R^qh'_*i'_*e_Z^{-1}i^{-1}\mathcal{F} =
R^q(h' \circ i')_*e_Z^{-1}i^{-1}\mathcal{F}'
$$
(first equality by
Lemma \ref{lemma-finite-pushforward-commutes-with-base-change}
and the second because
$Ri'_* = i'_*$ by
Proposition \ref{proposition-finite-higher-direct-image-zero})
to we see that we have $P(Z)$.
Finally, suppose $\xi'$ does not map to zero in the first summand.
We have
$$
(e')^{-1}j_*\mathcal{J} = j'_*e_V^{-1}\mathcal{J}
\quad\text{and}\quad
R^aj'_*e_V^{-1}\mathcal{J} = 0, \quad a = 1, \ldots, q - 1
$$
by $BC(f, n, q - 1)$ applied to the diagram
$$
\xymatrix{
X \ar[d]_f & Y' \ar[l] \ar[d]_{e'} & Y \ar[l]^{j'} \ar[d]^{e_V} \\
S & T' \ar[l] & V \ar[l]_j
}
$$
and the fact that $\mathcal{J}$ is injective.
By the relative Leray spectral sequence for $h' \circ j'$
(Cohomology on Sites, Lemma \ref{sites-cohomology-lemma-relative-Leray})
we deduce that
$$
R^qh'_*(e')^{-1}j_*\mathcal{J} =
R^qh'_*j'_*e_V^{-1}\mathcal{J}
\longrightarrow
R^q(h' \circ j')_* e_V^{-1}\mathcal{J}
$$
is injective. Thus $\xi$ maps to a nonzero element of
$(R^q(h' \circ j')_* e_V^{-1}\mathcal{J})_{\overline{x}}$.
Applying part (3) of Lemma \ref{lemma-base-change-q-injective}
to the injection $j^{-1}\mathcal{F}' \to \mathcal{J}$
we conclude that $P(V)$ holds.
\end{proof}
